% !TeX root = ../../tesis.tex
% §2.3.2 — movido de 2-MarcoTeorico/2-4-Teoria-Casos.tex l. 3-24 (#36)

\phantomsection\label{sec:teoria-latinoamericana}

Para enmarcar el caso de la \cdmx{} en un contexto de crítica urbanística
regional, es fundamental recurrir a la teoría sobre la ciudad en América Latina.
La obra compilatoria de \citet{padilla2007} proporciona el marco conceptual para
entender que la crisis de movilidad en la periferia es síntoma de un agotamiento
del modelo de planeación estatal y la adopción de enfoques que favorecieron el
capital inmobiliario y la privatización del espacio. La incapacidad de la
planeación urbana para gestionar la expansión territorial generada por los
grandes proyectos viales, como el Periférico, es el caldo de cultivo para la
emergencia de la ciudad neoliberal.

El escenario de un corredor anillar constituye, en este contexto teórico, un
caso de demostración metodológica: permite evaluar si la transición de un
enfoque radial ---propio de la era Uruchurtu--- hacia uno que incorpora
geometrías transversales reduce cuantitativamente los indicadores de
ineficiencia topológica que aquejaban a la ciudad fragmentada
\citep{franco2022}.
